%------------------------------------------------------------------------------%
\begin{exercises}

%------------------------------------------------------------------------------%
\begin{exercise}
Calcule as integrais.
\begin{mathlist}
  \item \int \sin^3(x)\cos^2(x)dx
  \item \int \sin^6(x)\cos^3(x)dx
  \item \int_0^{\frac{\pi}{2}} \sin^7(x) \cos^5(x) dx
  \item \int_0^{\frac{\pi}{2}} \sin^2(\frac{1}{3}x) dx
  \item \int_0^{\pi} cos^4(2t) dt
  \item \int t \sin^2(t)dt
  \item \int \cos(\theta) \cos^5(\sin(\theta)) d\theta
  \item \int \cos^2(x) \sin(2x)dx
  \item \int \tan^4(x) \sec^6(x) dx
  \item \int_0^{\frac{\pi}{3}} \tan^5(x) \sec^4(x)dx
  \item \int \sin(8x) \cos(5x)dx
  \item \int \frac{\sin(\phi)}{\cos^3(\phi)} d \phi
  \item \int \tan^5(x)dx
  \item \int_{\frac{\pi}{6}}^{\frac{\pi}{2}} \cot^2(x) dx
\end{mathlist}
\begin{answer}
\begin{mathlist}
  \item -\frac{1}{3} \cos^2(x)+\frac{1}{5} \cos^5(x)+C
  \item \frac{1}{7} \sin^7(x)-\frac{1}{5} \sin^5(x)+C
  \item \frac{1}{120}
  \item \pi+ \frac{3}{8} \sqrt{3}
  \item \frac{3}{8} \pi
  \item \frac{1}{4}t^2-\frac{1}{4} t \sin(2t)- \frac{1}{8} \cos(2t)+C
  \item \sin(\sin(\theta))- \frac{2}{3}\sin^5(\sin(\theta))+ \frac{1}{5} \sin^5(\sin(\theta))+C
  \item -\frac{1}{2}\cos^4(x)+C
  \item  \frac{1}{9} \tan^9(x)+ \frac{2}{7} \tan^7(x)+ \frac{1}{5} \tan^5(x)+C
  \item \frac{117}{8}
  \item -\frac{1}{6} \cos(3x)-\frac{1}{26} \cos(13x)+C
  \item \frac{1}{2} \tan^2(x)+C
  \item \frac{1}{4} \sec^4(x)- \tan^2(x) + ln \vert \sec(x) \vert+ C
  \item \frac{1}{2}(1-\ln 2)
\end{mathlist}
\end{answer}
\end{exercise}

%------------------------------------------------------------------------------%
\begin{exercise}
Uma partícula se move em linha reta com função velocidade
$v(t)= \sin(\omega t) \cos^2(\omega t)$.
Encontre a sua função posição $s=f(t)$ se $f(0)=0$.
\begin{answer}
  $s=\frac{(1-\cos^3(\omega t))}{3\omega}$
\end{answer}
\end{exercise}

%------------------------------------------------------------------------------%
\begin{exercise}
Determine se a $ \int_{0}^{\infty} \sin^2(x) dx$ integral é convergente ou divergente. Caso seja convergente, calcule o valor dessa integral.
\begin{answer}
divergente.
\end{answer}
\end{exercise}
\begin{exercise}
Nos exercícios seguintes encontre a área da região delimitada pelas curvas dadas e encontre o volume do sólido obtido pela rotação da região delimitada pelas curvas dadas em torno dos eixos especificados:
\begin{itemize}
    \item[a)] $y=\sin(x)$, $y=0$, $\frac{\pi}{2} \leq x \leq \pi$, em torno do eixo $x$.
    \item[b)] $y=\sin(x)$, $y=\cos(x)$, $0\leq x \leq \frac{\pi}{4} $, em torno de y=1.
\end{itemize}
\begin{answer}
\begin{itemize}
    \item[a)] $A=1$ e $V=\frac{\pi^2}{4}$
    \item[b)] $A=\sqrt{2}-1$ e $V=\pi(2\sqrt{2}- \frac{5}{2})$
\end{itemize}
\end{answer}
\end{exercise}

\begin{exercise}
Encontre o valor médio de $f(x)= \sin^2(x)\cos^3(x)$ no intervalo $[-\pi,\pi]$.
\begin{answer}
0
\end{answer}
\end{exercise}
\begin{exercise}
 Considere a região do primeiro quadrante delimitada pelos eixoss coordenadas e pela curva $f(x)=\cos(x)$, $0 \leq x \leq \frac{\pi}{2}$.
    \begin{itemize} 
        \item[a)] Encontre a área dessa região.
        \item[b)] Encontre o volume do sólido de revolução gerado pela rotação dessa região em torno do eixo $x$.
        \item[b)] Encontre o volume do sólido de revolução gerado pela rotação dessa região em torno do eixo $y$.
         \item[c)] Encontre o volume do sólido de revolução gerado pela rotação dessa região em torno do da reta $x=\frac{\pi}{2}$.
    \end{itemize}
\end{exercise}
%------------------------------------------------------------------------------%
%\begin{exercise}
%\begin{mathlist}[2]
%  \item
%  \item
%  \item
%\end{mathlist}
%\begin{answer}
%  \begin{alphalist}[3]
  %    \item
  %    \item
  %    \item
  %  \end{alphalist}
%\end{answer}
%\end{exercise}

\end{exercises}
%------------------------------------------------------------------------------%
